\documentclass[literature.tex]{subfiles}

\begin{document}
\prose{Báječná léta pod psa}{Michal Viewegh}
\teorieA{
epická próza, satirický autobiografický román (rodinná kronika s prvky grotesky a humoru)
}{}{
Hovorový, živý a sarkastický jazyk blízký čtenáři. Častá přímá řeč, dialogy, ironie a sebeironie. Mísí spisovný jazyk s vulgarismy, anglicismy a germanismy. Některé kapitoly mají formu divadelního scénáře nebo výpisků z deníku.
}{
Ironie, sarkasmus, sebeironie, opakování, pointy a dialogy s gradací.
}{
Přirovnání, metafora (rakev jako symbol zoufalství a úniku), ironie („báječná léta“ jako označení normalizace).
}{
Retrospektivní vyprávění v subjektivizované er-formě z pohledu Kvida s rámcovými dialogy s nakladatelem. Rychlý spád, humoristický a tragikomický tón. Prokládáno epizodami (babička Líba, sexuální dospívání, recitace). Dominuje ironický nadhled a sebeironie.
}{
„Víš, proč tam nejedeš, i když seš lepší?“
Kvidova otce známým způsobem píchlo pod prsní kostí.
„Dej pokoj,“ řekl co nejlhostejněji.
„Myslím to vážně,“ řekl Zvára. „Víš proč?“
„Nevím, šéfe. Nevysvětlil byste mi to?“ pokoušel se Kvidův otec o lehčí tón, ale náhlé napětí v kanceláři se tím nezmírnilo.
„Protože seš vůl!“
„Vidíš,“ řekl Kvidův otec malinko přiškrceným hlasem, „a já jsem si, hlupák, celou dobu myslel, že nejedu, protože nejsem v komunistické straně.“
}
\teorieB{
\wtem{Kvido}{hlavní postava a autorovo alter ego, nadprůměrně inteligentní, neohrabaný a tlustý kluk (později mladý muž), posedlý psaním a literaturou. Citlivý kronikář rodiny, zamilovaný do Jarušky, v dospělosti opouští VŠE kvůli psaní a rodině.}
\wtem{Otec}{nejrozporuplnější postava, inteligentní ekonom s červeným diplomem, neprůbojný introvert. Odmítá vstoupit do KSČ, ale kvůli rodině dělá kompromisy. Krátce „světák“ díky zahraničním cestám, pak se zhroutí a vyrábí si rakev v dílně.}
\wtem{Matka}{praktická, milující právnička a bývalá herečka. Bojí se psů (trauma z těhotenství), organizuje rodinné „záchranné akce“.}
\wtem{Jaruška}{Kvidova celoživotní láska, spolužačka, alergická na všechno. Z dětské kamarádky se stává manželkou a matkou Aničky.}
\wtem{Paco}{mladší bratr, průměrný, „kovbojský“ typ.}
\wtem{Babička}{ortodoxní vegetariánka, turistka a šetřivá cestovatelka, zdroj komických epizod (zeleninová dieta, strach z hliníku).}
\wtem{Šperk a Zvára}{typičtí představitelé režimu (stranický tajemník, pragmatický kolega).}
}{
Román je rodinnou kronikou od narození Kvida (na jevišti při absurdální hře Čekání na Godota) přes srpen 1968 až po sametovou revoluci a rané 90. léta. Rodina utíká z Prahy do Sázavy kvůli normalizaci. Kvido dospívá, recituje socialistickou poezii (pomáhá rodině k bytu), zamiluje se do Jarušky. Otec dělá morální kompromisy pro lepší bydlení a práci (vstup do VUML, nákup psa od straníka, fotbal), krátce „procestuje“ svět, ale po návštěvě u chartisty Pavla Kohouta je degradován na vrátného a psychicky se zhroutí – zavírá se do dílny a vyrábí si ozdobenou rakev. Kvido opouští ekonomku, píše, s Jaruškou mají dceru Aničku (na naléhání matky, aby „zachránili“ otce). Paralelně probíhají dialogy dospělého Kvida s nakladatelem o cenzuře. Epilog shrnuje osudy po roce 1989 formou absurdního výčtu.
}{
17 číslovaných kapitol + epilog. Rámcová kompozice (Kvido předčítá román nakladateli) s retrospektivním vyprávěním. Dvě linie: hlavní rodinný příběh a metarozhovory o vydání. Prostor: Praha → Sázava (Posázaví). Čas: 1962 (narození) – 1968–1989 (normalizace) – 1991 (epilog).
}{
Ironický a sebeironický pohled na život v normalizačním Československu. Ukazuje, jak režim deformuje charaktery (kompromisy, strach, psychické zhroucení), ale zároveň oslavuje rodinnou soudržnost a humor jako obranu. Motivy: maloměšťáctví, cenzura, dospívání, láska, absurdita totality (\texthl{„báječná léta pod psa“} jako oxymóron). Viewegh říká: „převážně veselé ohlédnutí za převážně neveselou dobou“.
}\newpage
\historie{
Dílo zachycuje normalizaci po roce 1968 (sovětská okupace, perzekuce, cenzura) až po sametovou revoluci 1989. Kritizuje morální ústupky, VUML, kádrování a absurditu režimu.
}{}{}{}{}{
Michal Viewegh (* 1962 Praha) je český spisovatel a publicista. Vystudoval VŠE (nedokončil) a bohemistiku/pedagogiku na FF UK. Pracoval jako vrátný, učitel, redaktor. Debutoval právě touto knihou. Jeho tvorba často čerpá z autobiografických prvků a satiricky zobrazuje českou společnost.
}{
Mezi nejznámější díla patří \textsc{Výchova dívek v Čechách}, \textsc{Účastníci zájezdu}, \textsc{Román pro ženy}, \textsc{Můj život po životě} nebo \textsc{Děravé paměti}.
}
\kritika{}{}{
Témata morálních kompromisů, vlivu politiky na rodinu a humor jako obrany zůstávají relevantní. Kniha (a její filmová adaptace) patří k nejčtenějším českým románům 90. let.
}
\nazor{
Vieweghova prvotina mě bavila od první stránky. Skvěle kombinuje sarkastický humor, ironii a tragikomiku – otec vyrábějící si rakev je geniální obraz zoufalství z totality. Nejsilnější je kontrast mezi dětskou naivitou Kvida (recitace socialistické poezie) a dospělým nadhledem. Dialogy s nakladatelem o cenzuře jsou trefné. Kniha není jen „vtipná“, ale i smutná a lidská. Rozhodně doporučuji – je to jedno z nejlepších ohlédnutí za „báječnými“ normalizačními lety.
}
\explain{
\textbl{Satira --} literární žánr zesměšňující společenské nedostatky;
\textbl{Groteska --} spojení komického a tragického, často absurdního;
\textbl{Autobiografický román --} dílo s prvky autorova vlastního života.
}
\wpage
\end{document}